\subsection{Direct probes for a sharp five-copy certificate}
\label{sec:fivecopy-direct-probes}

Let $W=(\mathbb C^2)^{\otimes5}$, let $U_\sigma$ permute its five
copies, and let $G$ be any element of the rational local commutant
$\mathcal A_5$. On seven sites put
\[
 P=\frac1{120}\sum_{\sigma\in S_5}U_\sigma^{\otimes7},\qquad
 X(G)=PG^{\otimes7}P.
\]
These linear probes determine identities on the moment space; they need not themselves be positive moments. For $j=0,1,2$, let
$H_j=S_{5-j}\times S_j$, let $\phi_j$ be the local $j$-copy spin-flip transpose,
and put $\Phi_j=\phi_j^{\otimes7}$.

\begin{lemma}[direct contracted dual probes; DP5-1]
\label{lem:fivecopy-direct-probes}
Suppose $Z$ is a site-permutation-invariant commutant operator supported
on the global $H_j$-invariant subspace. Choose left-coset representatives
$a\in H_j\backslash S_5$ and right-coset representatives
$b\in S_5/H_j$. Then
\[
 \operatorname{tr}[Z\Phi_j(X(G))]
 =\binom5j^{-2}\sum_{a,b}
       \operatorname{tr}\left[Z\,\phi_j(U_aGU_b)^{\otimes7}\right].
\]
There are $1,25,100$ terms for $j=0,1,2$, respectively.
The traces on the right are evaluated directly in the chosen Schur models;
no full gamma map is needed. Write
$a_\sigma=\operatorname{tr}(U_\sigma G)$, and let $\mathcal T$ be the
ten transpositions in $S_5$. The sharp target on this probe is
\[
 q(G)=\frac1{1200}\sum_{\sigma\in S_5,\,t\in\mathcal T}
       a_{t\sigma}^3a_\sigma^4
       -\frac{121}{917\cdot120}\sum_{\sigma\in S_5}a_\sigma^7.
\]
Equality of this polynomial with the sum of the contracted dual probes
for every $G\in\mathcal A_5$ is equivalent to the sharp dual identity
on the real, site-averaged bosonic moment space.
\end{lemma}
\begin{proof}
The spin-flip transpose is self-adjoint for the physical trace pairing.
Expanding both copies of $P$ gives $120^2$ terms. Both $Z$ and $\Phi_j(Z)$
have left and right $H_j$-invariant support: partial transpose exchanges
left and right permutations of the flipped copies and preserves the
unflipped ones. The terms therefore group into $\binom5j^2$ equal-weight
coset pairs. Moving $\Phi_j$ back to the probe gives the first formula.
The same proof applies with any numbers of copies and sites in place of $5$ and $7$.

On the bosonic subspace, the purity operator may be averaged over all
ten choices of the swapped global-copy pair. This average commutes with
$P$, so cyclicity removes one of the two projectors in its trace against
$X(G)$. All $35$ choices of the three swapped sites have the same trace
against an identical-site product $G^{\otimes7}$, yielding the formula
for $q$. Finally the span of the powers $G^{\otimes7}$ is the symmetric
seventh tensor power of $\mathcal A_5$, by polarisation. The trace pairing
on this site-symmetric algebra is nondegenerate. This proves the last claim.
Probes over $\mathbb Q(i)$ are also allowed: real and imaginary parts of
$G=A+iB$ evaluations are rational combinations of the eight real probes
$A+tB$, $t=0,\ldots,7$, by degree-seven interpolation.
\end{proof}

\begin{proposition}[the critical space and finite probes; DP5-1]
\label{prop:fivecopy-finite-probes}
Let $\mathcal M_5$ be the real span of the normalised five-copy moments
of pure seven-qubit states after local-unitary twirling, site averaging
and averaging with the complex conjugate state (the conjugation-even
part). Then $\dim\mathcal M_5=3285$. The \emph{critical space}
$\mathcal C\subset\mathcal M_5^*$ of functionals vanishing at the
icosahedral moment $X_0$ and at its derivative $X_0'$ along the
normalised family $\psi_x$ of~\eqref{eq:ico} at $x_0=37/524$ has
dimension $3283$. If functionals in $\mathcal C$ have an evaluation
matrix of rank $3283$ on some finite set of probes, they span
$\mathcal C$, and evaluation on those probes is injective on
$\mathcal C$.
\end{proposition}
\begin{proof}
By the last part of Lemma~\ref{lem:fivecopy-direct-probes}, $\mathcal M_5$
is dual to the site-symmetric probe polynomials, that is, to the permutation-symmetric local-unitary
invariants of bidegree $(5,5)$ for $N=7$ sites and local dimension $r=2$. Their Schur--Weyl character sum
(\eqref{eq:fqe-invariant-dimension} and~\eqref{eq:fqe-invariant-even} of
Section~\ref{sec:fqe-invariant-dimensions}) gives dimension $5844$, and the trace of the conjugation swap is
$726$ (DP5-1), so the conjugation-even part has dimension
$(5844+726)/2=3285$. This is the dimension of the linear span $\mathcal M_5$; the normalised moments
lie in its affine hyperplane $\nu=1$. The normalisation functional $\nu$ has
$\nu(X_0)=1$ and $\nu(X_0')=0$, and DP5-1 checks that the weight-two Pauli enumerator
$A_2$ has derivative $96/131\ne0$ there, so $X_0'\ne0$ and $X_0,X_0'$ are independent. Hence
$\dim\mathcal C=3283$. A rank-$3283$ evaluation matrix of elements of
$\mathcal C$ gives both remaining assertions.
\end{proof}

The criterion is applied to the twelve-face baseline in
Section~\ref{sec:fivecopy-positive-repair}, where rank and repair bounds
are verified for the rational data.
